\documentclass[10pt,a4paper]{amsart}
\usepackage{amsmath,amssymb,amsthm,graphicx}
% For commutative diagrams we use Xy-pic
\usepackage[all]{xy}

% We use 2cell for 2-commutative diagrams.
\xyoption{2cell}
% Theorem environments.
%
\theoremstyle{plain}
\newtheorem{theorem}[subsection]{Theorem}
\newtheorem{proposition}[subsection]{Proposition}
\newtheorem{lemma}[subsection]{Lemma}

\theoremstyle{definition}
\newtheorem{definition}[subsection]{Definition}
\newtheorem{example}[subsection]{Example}
\newtheorem{exercise}[subsection]{Exercise}
\newtheorem{situation}[subsection]{Situation}

\theoremstyle{remark}
\newtheorem{remark}[subsection]{Remark}
\newtheorem{remarks}[subsection]{Remarks}

\numberwithin{equation}{subsection}

% Macros
%
\def\lim{\mathop{\mathrm{lim}}\nolimits}
\def\colim{\mathop{\mathrm{colim}}\nolimits}
\def\Spec{\mathop{\mathrm{Spec}}}
\def\Hom{\mathop{\mathrm{Hom}}\nolimits}
\def\Ext{\mathop{\mathrm{Ext}}\nolimits}
\def\SheafHom{\mathop{\mathcal{H}\!\mathit{om}}\nolimits}
\def\SheafExt{\mathop{\mathcal{E}\!\mathit{xt}}\nolimits}
\def\Sch{\mathit{Sch}}
\def\Mor{\mathop{\mathrm{Mor}}\nolimits}
\def\Ob{\mathop{\mathrm{Ob}}\nolimits}
\def\Sh{\mathop{\mathit{Sh}}\nolimits}
\def\NL{\mathop{N\!L}\nolimits}
\def\CH{\mathop{\mathrm{CH}}\nolimits}
\def\proetale{{pro\text{-}\acute{e}tale}}
\def\etale{{\acute{e}tale}}
\def\QCoh{\mathit{QCoh}}
\def\Ker{\mathop{\mathrm{Ker}}}
\def\Im{\mathop{\mathrm{Im}}}
\def\Coker{\mathop{\mathrm{Coker}}}
\def\Coim{\mathop{\mathrm{Coim}}}

% Boxtimes
%
\DeclareMathSymbol{\boxtimes}{\mathbin}{AMSa}{"02}

%
% Macros for moduli stacks/spaces
%
\def\QCohstack{\mathcal{QC}\!\mathit{oh}}
\def\Cohstack{\mathcal{C}\!\mathit{oh}}
\def\Spacesstack{\mathcal{S}\!\mathit{paces}}
\def\Quotfunctor{\mathrm{Quot}}
\def\Hilbfunctor{\mathrm{Hilb}}
\def\Curvesstack{\mathcal{C}\!\mathit{urves}}
\def\Polarizedstack{\mathcal{P}\!\mathit{olarized}}
\def\Complexesstack{\mathcal{C}\!\mathit{omplexes}}
% \Pic is the operator that assigns to X its picard group, usage \Pic(X)
% \Picardstack_{X/B} denotes the Picard stack of X over B
% \Picardfunctor_{X/B} denotes the Picard functor of X over B
\def\Pic{\mathop{\mathrm{Pic}}\nolimits}
\def\Picardstack{\mathcal{P}\!\mathit{ic}}
\def\Picardfunctor{\mathrm{Pic}}
\def\Deformationcategory{\mathcal{D}\!\mathit{ef}}
\newenvironment{reference}{}{}
\newenvironment{slogan}{}{}
\newenvironment{history}{}{}
\usepackage{cleveref}
\crefname{originaldefinition}{original definition}{original definitions}
\Crefname{originaldefinition}{Original definition}{Original definitions}
\makeatletter
\newlabel{divisors-definition-strict-transform}{{080D}{source}{}{sourcepointer-1}{}}
\newlabel{divisors-definition-strict-transform@cref}{{[originaldefinition][0][]080D}{[1][0][]source}{}{}{}}
\makeatother
\crefname{originallemma}{original lemma}{original lemmas}
\Crefname{originallemma}{Original lemma}{Original lemmas}
\makeatletter
\newlabel{divisors-lemma-blowing-up-projective}{{02NS}{source}{}{sourcepointer-2}{}}
\newlabel{divisors-lemma-blowing-up-projective@cref}{{[originallemma][0][]02NS}{[1][0][]source}{}{}{}}
\makeatother
\crefname{originallemma}{original lemma}{original lemmas}
\Crefname{originallemma}{Original lemma}{Original lemmas}
\makeatletter
\newlabel{chow-lemma-proper-pushforward-rational-equivalence}{{02S2}{source}{}{sourcepointer-3}{}}
\newlabel{chow-lemma-proper-pushforward-rational-equivalence@cref}{{[originallemma][0][]02S2}{[1][0][]source}{}{}{}}
\makeatother
\crefname{originallemma}{original lemma}{original lemmas}
\Crefname{originallemma}{Original lemma}{Original lemmas}
\makeatletter
\newlabel{chow-lemma-factors}{{02TI}{source}{}{sourcepointer-4}{}}
\newlabel{chow-lemma-factors@cref}{{[originallemma][0][]02TI}{[1][0][]source}{}{}{}}
\makeatother
\crefname{originaldefinition}{original definition}{original definitions}
\Crefname{originaldefinition}{Original definition}{Original definitions}
\makeatletter
\newlabel{groupoids-definition-principal-homogeneous-space}{{049A}{source}{}{sourcepointer-5}{}}
\newlabel{groupoids-definition-principal-homogeneous-space@cref}{{[originaldefinition][0][]049A}{[1][0][]source}{}{}{}}
\makeatother
\crefname{originallemma}{original lemma}{original lemmas}
\Crefname{originallemma}{Original lemma}{Original lemmas}
\makeatletter
\newlabel{descent-lemma-affine}{{0245}{source}{}{sourcepointer-6}{}}
\newlabel{descent-lemma-affine@cref}{{[originallemma][0][]0245}{[1][0][]source}{}{}{}}
\makeatother
\crefname{originallemma}{original lemma}{original lemmas}
\Crefname{originallemma}{Original lemma}{Original lemmas}
\makeatletter
\newlabel{more-algebra-lemma-regular-local-UFD}{{0AG0}{source}{}{sourcepointer-7}{}}
\newlabel{more-algebra-lemma-regular-local-UFD@cref}{{[originallemma][0][]0AG0}{[1][0][]source}{}{}{}}
\makeatother
\title{Projective schemes admit non-effective etale descent data}
\author{The Stacks Project Authors}
\date{stacks-project:examples.tex:local-snapshot-20261011; GNU FDL 1.2 or later}
\begin{document}
\maketitle
\noindent\textit{Editorial note.} The earlier non-locality-of-projectivity example is retained as uncounted proof support. This task counts the stronger non-effectivity of a projective descent datum, not the earlier example as a second task.\par\medskip
\noindent\textit{Allowed standard prerequisites (editorial).}\begin{itemize}
\item Explicit prerequisite boundary. Blowups of smooth centres and their strict transforms, finite-group torsors and invariant quotients, affine descent, regular-local factoriality, rational equivalence and degree of line bundles on proper cycles.
\end{itemize}\medskip
\ifdefined\hypertarget\hypertarget{sourcepointer-1}{}\fi
\noindent\textit{Source reference.} Original definition 080D: focused-seeds/AGP/sources/stacks-snapshot-20261011/tags.txt, source lines 10178--10178. Stable Stacks tag for an expressly allowed foundational or notational input; the counted construction remains the two concrete core obligations. The retained local core is not replaced by this pointer. See standard\_prerequisites.\par\medskip
\ifdefined\hypertarget\hypertarget{sourcepointer-2}{}\fi
\noindent\textit{Source reference.} Original lemma 02NS: focused-seeds/AGP/sources/stacks-snapshot-20261011/tags.txt, source lines 3491--3491. Stable Stacks tag for an expressly allowed foundational or notational input; the counted construction remains the two concrete core obligations. The retained local core is not replaced by this pointer. See standard\_prerequisites.\par\medskip
\ifdefined\hypertarget\hypertarget{sourcepointer-3}{}\fi
\noindent\textit{Source reference.} Original lemma 02S2: focused-seeds/AGP/sources/stacks-snapshot-20261011/tags.txt, source lines 3645--3645. Stable Stacks tag for an expressly allowed foundational or notational input; the counted construction remains the two concrete core obligations. The retained local core is not replaced by this pointer. See standard\_prerequisites.\par\medskip
\ifdefined\hypertarget\hypertarget{sourcepointer-4}{}\fi
\noindent\textit{Source reference.} Original lemma 02TI: focused-seeds/AGP/sources/stacks-snapshot-20261011/tags.txt, source lines 3697--3697. Stable Stacks tag for an expressly allowed foundational or notational input; the counted construction remains the two concrete core obligations. The retained local core is not replaced by this pointer. See standard\_prerequisites.\par\medskip
\ifdefined\hypertarget\hypertarget{sourcepointer-5}{}\fi
\noindent\textit{Source reference.} Original definition 049A: focused-seeds/AGP/sources/stacks-snapshot-20261011/tags.txt, source lines 5573--5573. Stable Stacks tag for an expressly allowed foundational or notational input; the counted construction remains the two concrete core obligations. The retained local core is not replaced by this pointer. See standard\_prerequisites.\par\medskip
\ifdefined\hypertarget\hypertarget{sourcepointer-6}{}\fi
\noindent\textit{Source reference.} Original lemma 0245: focused-seeds/AGP/sources/stacks-snapshot-20261011/tags.txt, source lines 2766--2766. Stable Stacks tag for an expressly allowed foundational or notational input; the counted construction remains the two concrete core obligations. The retained local core is not replaced by this pointer. See standard\_prerequisites.\par\medskip
\ifdefined\hypertarget\hypertarget{sourcepointer-7}{}\fi
\noindent\textit{Source reference.} Original lemma 0AG0: focused-seeds/AGP/sources/stacks-snapshot-20261011/tags.txt, source lines 13178--13178. Stable Stacks tag for an expressly allowed foundational or notational input; the counted construction remains the two concrete core obligations. The retained local core is not replaced by this pointer. See standard\_prerequisites.\par\medskip
\begin{lemma}
\label{lemma-non-effective-descent-projective}
There is an etale covering $X\to S$ of schemes and a descent datum
$(V/X,\varphi)$ relative to $X\to S$ such that 
$V\to X$ is projective,
but the descent datum is not effective in the category of schemes.
\end{lemma}

\begin{lemma}
\label{lemma-non-descending-property-projective}
The properties
\begin{enumerate}
\item[] $\mathcal{P}(f) =$``$f$ is projective'', and
\item[] $\mathcal{P}(f) =$``$f$ is quasi-projective''
\end{enumerate}
are not Zariski local on the base. A fortiori, they are not fpqc local
on the base.
\end{lemma}

\begin{proof}
Following Hironaka \cite[Example B.3.4.1]{H},
we define a proper morphism of smooth complex 3-folds
$f:V_Y\to Y$
which is Zariski-locally projective, but not projective. Since $f$ is proper
and not projective, it is also not quasi-projective.

\medskip\noindent
Let $Y$ be projective 3-space over the complex numbers ${\mathbf C}$.
Let $C$ and $D$ be smooth conics in $Y$ such
that the closed subscheme $C\cap D$ is reduced and consists
of two complex points $P$ and $Q$. (For example,
let $C=\{ [x,y,z,w]: xy=z^2, w=0\}$, $D=\{ [x,y,z,w]:
xy=w^2, z=0\}$, $P=[1,0,0,0]$,
and $Q=[0,1,0,0]$.) 
On $Y-Q$, first blow up the curve $C$, and then blow
up the strict transform of the curve $D$ (Divisors, Definition 
\ref{divisors-definition-strict-transform}). On $Y-P$, first blow up
the curve $D$, and then blow up the strict transform of the curve
$C$. Over $Y-P-Q$, the two varieties we have constructed are canonically
isomorphic, and so we can glue them over $Y-P-Q$. The result
is a smooth proper 3-fold $V_Y$ over ${\mathbf C}$. The morphism
$f:V_Y\to Y$ is proper and Zariski-locally projective (since
it is a blowup over $Y-P$ and over $Y-Q$), by Divisors,
Lemma \ref{divisors-lemma-blowing-up-projective}. We will show that
$V_Y$ is not projective over ${\mathbf C}$. That will imply that
$f$ is not projective.

\medskip\noindent
To do this, let $L$ be the inverse image in $V_Y$ of a complex point
of $C-P-Q$, and $M$ the inverse image of a complex point of $D-P-Q$.
Then $L$ and $M$ are isomorphic to the projective line
${\mathbf P}^1_{{\mathbf C}}$. 
Next, let $E$ be the inverse image in $V_Y$ of $C\cup D\subset Y$ in $V_Y$;
thus $E\rightarrow C\cup D$ is a proper morphism, with fibers
isomorphic to ${\mathbf P}^1$ over $(C\cup D)-\{P,Q\}$.
The inverse
image of $P$ in $E$ is a union of two lines $L_0$ and $M_0$, and we have
rational equivalences of cycles $L\sim L_0+M_0$ and $M\sim M_0$ on $E$
(using that $C$ and $D$ are isomorphic to ${\mathbf P}^1$).
Note the asymmetry resulting from the order in which we blew
up the two curves. Near $Q$, the opposite happens. So the inverse image
of $Q$ is the union of two lines $L_0'$ and $M_0'$, and we have
rational equivalences $L\sim L_0'$ and $M\sim L_0'+M_0'$ on $E$.
Combining these equivalences, we find that $L_0+M_0'\sim 0$
on $E$ and hence on $V_Y$. If $V_Y$ were projective over ${\mathbf C}$,
it would have
an ample line bundle $H$, which would have degree $> 0$ on all curves
in $V_Y$. In particular $H$ would have positive degree on $L_0+M_0'$,
contradicting that the degree of a line bundle is well-defined
on 1-cycles modulo rational equivalence on a proper scheme
over a field (Chow Homology,
Lemma \ref{chow-lemma-proper-pushforward-rational-equivalence}
and Lemma \ref{chow-lemma-factors}).
So $V_Y$ is not projective over ${\mathbf C}$.
\end{proof}

\begin{proof}
We imitate Hironaka's example of a smooth separated complex
algebraic space of dimension 3
which is not a scheme \cite[Example B.3.4.2]{H}.

\medskip\noindent
Consider the action of the group $G = \mathbf{Z}/2 = \{1, g\}$
on projective 3-space
$\mathbf{P}^3$ over the complex numbers by
$$
g[x,y,z,w] = [y,x,w,z].
$$
The action is free outside the two disjoint lines
$L_1=\{ [x,x,z,z]\}$ and $L_2=\{ [x,-x,z,-z]\}$ in
${\mathbf P}^3$. Let $Y={\mathbf P}^3-(L_1\cup L_2)$. There is a
smooth quasi-projective scheme $S=Y/G$ over ${\mathbf C}$ such that
$Y\to S$ is a $G$-torsor (Groupoids,
Definition \ref{groupoids-definition-principal-homogeneous-space}).
Explicitly, we can define $S$ as the image of the open subset $Y$
in ${\mathbf P}^3$ under the morphism
\begin{align*}
{\mathbf P}^3 & \to \text{Proj } {\mathbf C}[x,y,z,w]^G\\
   & = \text{Proj } {\mathbf C}[u_0,u_1,v_0,v_1,v_2]/(v_0v_1=v_2^2),
\end{align*}
where $u_0=x+y$, $u_1=z+w$, $v_0=(x-y)^2$, $v_1=(z-w)^2$,
and $v_2=(x-y)(z-w)$, and the ring is graded with $u_0,u_1$
in degree 1 and $v_0,v_1,v_2$ in degree 2.

\medskip\noindent
Let $C=\{ [x,y,z,w]: xy=z^2, w=0\}$ and $D=\{ [x,y,z,w]:
xy=w^2, z=0\}$. These are smooth conic curves in ${\mathbf P}^3$, contained
in the $G$-invariant open subset $Y$, with $g(C)=D$. Also,
$C\cap D$ consists of the two points $P:=[1,0,0,0]$
and $Q:=[0,1,0,0]$, and these two points are switched by the action
of $G$. 

\medskip\noindent
Let $V_Y\to Y$ be the scheme which over $Y-P$
is defined by blowing up $D$ and then the strict transform
of $C$, and over $Y-Q$ is defined by blowing up $C$ and then
the strict transform of $D$. (This is the same construction
as in the proof of Lemma \ref{lemma-non-descending-property-projective},
except that $Y$ here denotes an open subset of ${\mathbf P}^3$
rather than all of ${\mathbf P}^3$.)
Then the action of $G$ on $Y$ lifts to an action of $G$ on $V_Y$,
which switches the inverse images of $Y-P$ and $Y-Q$. This action
of $G$ on $V_Y$ gives a descent datum $(V_Y/Y,\varphi_Y)$
on $V_Y$ relative to the $G$-torsor
$Y\to S$. The morphism $V_Y\to Y$
is proper but not projective, as shown in the proof
of Lemma \ref{lemma-non-descending-property-projective}.

\medskip\noindent
Let $X$ be the disjoint union of the open subsets $Y-P$ and $Y-Q$;
then we have surjective etale morphisms $X\to Y\to S$.
Let $V$ be the pullback of $V_Y\to Y$ to $X$; then the morphism
$V\to X$ is projective, since $V_Y\to Y$ is a blowup over each of the open
subsets $Y-P$ and $Y-Q$. Moreover, the descent datum $(V_Y/Y,\varphi_Y)$
pulls back to a descent datum $(V/X,\varphi)$ relative to the
etale covering $X\to S$.

\medskip\noindent
Suppose that this descent datum is effective in the category
of schemes. That is, there is a scheme $U\to S$
which pulls back to the morphism $V\to X$ together
with its descent datum. Then $U$ would be the quotient
of $V_Y$ by its $G$-action.
$$
\xymatrix{
V \ar[r]\ar[d]&  X\ar[d] \\
V_Y \ar[r]\ar[d]& Y\ar[d] \\
U \ar[r]& S
}
$$

\medskip\noindent
Let $E$ be the inverse image of $C\cup D\subset Y$ in $V_Y$;
thus $E\rightarrow C\cup D$ is a proper morphism, with fibers
isomorphic to ${\mathbf P}^1$ over $(C\cup D)-\{P,Q\}$.
The inverse image of $P$ in $E$ is a union of two lines $L_0$
and $M_0$. It follows that the inverse image of $Q=g(P)$ in $E$
is the union of two lines $L_0'=g(M_0)$ and $M_0'=g(L_0)$.
As shown in the proof
of Lemma \ref{lemma-non-descending-property-projective},
we have a rational equivalence $L_0+M_0'=L_0+g(L_0)\sim 0$ on $E$.

\medskip\noindent
By descent of closed subschemes, there is a curve $L_1\subset U$
(isomorphic to ${\mathbf P}^1$)
whose inverse image in $V_Y$ is $L_0\cup g(L_0)$. (Use Descent, Lemma
\ref{descent-lemma-affine}, noting that a closed immersion is an affine
morphism.)
Let $R$ be a complex point of $L_1$. Since
we assumed that $U$ is a scheme, we can choose a function
$f$ in the local ring $O_{U,R}$ that vanishes at $R$ but not
on the whole curve $L_1$. Let $D_{\text{loc}}$ be an irreducible component
of the closed subset $\{f = 0\}$ in $\Spec O_{U,R}$; then
$D_{\text{loc}}$ has codimension 1.
The closure of $D_{\text{loc}}$ in $U$ is an irreducible divisor $D_U$
in $U$ which contains the point $R$ but not the whole curve $L_1$.
The inverse image of $D_U$ in $V_Y$ is an effective divisor $D$
which intersects $L_0\cup g(L_0)$ but does not contain either
curve $L_0$ or $g(L_0)$.

\medskip\noindent
Since the complex 3-fold $V_Y$ is smooth, $O(D)$ is a line
bundle on $V_Y$. We use here that a regular local ring is factorial,
or in other words is a UFD, see
More on Algebra, Lemma \ref{more-algebra-lemma-regular-local-UFD}.
The restriction of $O(D)$ to the proper surface
$E\subset V_Y$ is a line bundle which has positive degree on the 1-cycle
$L_0+g(L_0)$, by our information on $D$. Since
$L_0+g(L_0)\sim 0$ on $E$, this contradicts 
that the degree of a line bundle is well-defined
on 1-cycles modulo rational equivalence on a proper scheme
over a field (Chow Homology,
Lemma \ref{chow-lemma-proper-pushforward-rational-equivalence}
and Lemma \ref{chow-lemma-factors}). Therefore the descent datum
$(V/X,\varphi)$ is in fact not effective; that is, $U$ does not exist
as a scheme.
\end{proof}

\bibliographystyle{plain}
\bibliography{my}
\end{document}
